\section{The Deterministic Service-Code Frontier}
Consider three dates. The root tier is fixed at zero. At the intermediate date a branch $j$ is observed with probability $\pi_j>0$, where $\sum_j\pi_j=1$. All branches then enter the same terminal public vertex. Payments and discounting equal one. Intermediate and terminal tiers are in $[0,1]$. Branch $j$ has cap $0<b_1<\cdots<b_k<1$, and the root promise is $\bar b=\sum_j\pi_jb_j$. At the intermediate date the reward is $-h_j(y)$, where $h_j$ is convex, nondecreasing, and $h_j(0)=0$. The common terminal reward $f(c)$ is continuous, concave, and strictly increasing on $[0,1]$. The terminal cap is one. The terminal decoder observes only its public vertex and its memory symbol, not the inherited intermediate tier.

\begin{theorem}[Heterogeneous renewal frontier]\label{thm:frontier}
The unique full-information policy is $y_j=0$, $c_j=b_j$, with value $J=\sum_j\pi_j f(b_j)$. Exact implementation requires $k$ symbols. With at most $m$ deterministic terminal symbols, an optimal encoding has contiguous nonempty cells in cap order. For a cell $i,\ldots,j$, its terminal tier is $b_i$ and its loss is
\begin{equation}\label{eq:general-cell}
 C(i,j)=\sum_{h=i}^j\pi_h\left[f(b_h)-f(b_i)+h_h(b_h-b_i)\right].
\end{equation}
The exact minimum loss is $D_m=D(\min(m,k),k)$, where
\begin{equation}\label{eq:segmentation}
 D(s,j)=\min_{s-1\leq i<j}\{D(s-1,i)+C(i+1,j)\},\quad D(0,0)=0,
\end{equation}
and all other infeasible entries equal $+\infty$. With a precomputed cell-cost table this requires $O(mk^2)$ arithmetic operations. For quadratic $f$ and quadratic $h_j$, prefix sums of the probability-weighted coefficients compute each cell cost in constant arithmetic time.
\end{theorem}
\begin{proof}
The root promise is the mean of the branch caps, so every branch cap must bind. Hence $y_j=b_j-c_j$ and $0\leq c_j\leq b_j$. The branch reward $f(c)-h_j(b_j-c)$ is strictly increasing on that interval. Full information therefore uniquely selects $c_j=b_j$; different terminal actions require different symbols.

Within a memory cell all terminal actions coincide and cannot exceed its smallest cap. Monotonicity selects that cap and subtraction from the full-information reward gives \eqref{eq:general-cell}. Fix any collection of cell minima. Every branch strictly prefers the largest eligible selected minimum, irrespective of its probability and intermediate cost. Reassigning branches in this way keeps each selected minimum in its own nonempty cell and produces contiguous cells without reducing value. Splitting at the final cell gives \eqref{eq:segmentation}. Expanding the quadratic terms gives the stated prefix-sum implementation.
\end{proof}

This strengthens and includes the quadratic equal-probability frontier. With $f(c)=2c-c^2/2$, $h_j(y)=y^2/2$, and $\pi_j=1/k$, \eqref{eq:general-cell} becomes
\begin{equation}\label{eq:old-cell}
 C(i,j)=\frac{2-b_i}{k}\sum_{h=i}^j(b_h-b_i).
\end{equation}
For $k=8$ and $b_j=j/9$, one, two, and three symbols lose $119/162$, $5/18$, and $49/324$, respectively. Two writable bits permit four symbols, not three. The segmentation recurrence is standard; the contract-specific cell loss, ordered-cell sufficiency, and exact memory interpretation are the substantive statements.

The earlier nonsaturated two-branch example also remains useful. With branch caps $1/4$ and $1$, root promise $1/2$, and the same quadratic rewards, the full optimum uses terminal tiers $1/4$ and $3/4$ and has value $27/32$. An optimized public-only table uses common terminal tier $1/4$ and intermediate tiers $0,1/2$, giving $13/32$. The adaptation gain is $7/16$ relative to the optimized restriction. At the full optimum both inherited intermediate tiers are zero, so making that tier observable does not evade this example's exact-memory lower bound. By contrast, the entire restricted-memory frontier above depends on not allowing a branch-specific intermediate tier to serve as an uncharged code.

